\section{Sezione critica e Peterson}
Come già menzionato, l'esecuzione di processi concorrentemente o in parallelo su dati condivisi può portare a problemi riguardanti l'integrità di questi
ultimi; la soluzione proposta nei capitoli precedenti implementava un buffer di dimensione fissa, ma sarebbe possibile anche utilizzare un contatore
che si incrementa ad ogni nuovo elemento inserito nel buffer e decrementa quando ne vengono rimossi.\par
Tuttavia, se il counter è inserito in più file ed è condiviso, è possibile che si presentino comportamenti indesiderati, come per esempio l'aumento e
diminuzione della variabile allo stesso tempo, riportando valori errati. Si ritiene dunque necessario fornire una garanzia che un solo processo alla
volta possa modificare i dati condivisi, evitando che questi dipendano dall'ordine degli accessi, quindi dalle race condition. La forma presa da queste
soluzioni ha il nome di \textbf{sincronizzazione e coordinamento tra processi}.\newline

\noindent Supponiamo un sistema composto da un totale di $n$ processi, ognuno dei quali con un proprio segmento di codice. Definiamo
\textbf{sezione critica} la parte del segmento in cui i processi possono accedere e modificare dati condivisi; la soluzione relativa alla gestione di
questo spazio consiste nel progettare un meccanismo per regolarne gli accessi. Esistono più sezioni di codice: \begin{itemize}
    \item \textbf{Sezione di ingresso}: Implementa il meccanismo di sincronizzazione.
    \item \textbf{Sezione di uscita}: Segue la sezione di ingresso.
    \item \textbf{Sezione non critica}: La parte restante del codice.
\end{itemize}

\noindent Il protocollo ideato per il software dovrà inoltre rispettare i seguenti tre requisiti: \begin{enumerate}
    \item \textbf{Mutua esclusione}: Quando un processo è in esecuzione all'interno della sua sezione critica, nessun altro processo può accedere alla
    propria.
    \item \textbf{Progresso}: Se nessun processo è in esecuzione nella propria sezione critica e qualcuno di loro desidera accedervi, solo i processi al
    di fuori delle relative sezioni critiche possono contribuire alla scelta di quello che potrà accedervi per primo. Questa non è una scelta
    infinitamente rimandabile.
    \item \textbf{Attesa limitata}: Se un processo ha fatto richiesta di entrare nella sua sezione critica, si pone un limite al numero di volte in cui
    gli altri processi possono accedere alle proprie prima che si accordi la richiesta del primo. In tal modo, tutti i processi potranno eventualmente
    accedere alla sezione.
\end{enumerate}

\noindent La soluzione del problema in ambiente single-core è banalmente ottenuta impedendo il verificarsi di interruzioni durante la modifica di dati
condivisi, quindi fornendo algoritmi senza prelazione. Tuttavia ciò non si può dire per sistemi multi-core, perché disabilitare gli interrupt potrebbe
richiedere molto tempo, impattando anche l'efficienza del sistema.\par
Le due soluzioni proposte per il problema della sezione critica in ambiente multi-core sono: \begin{itemize}
    \item \textbf{Kernel con diritto di prelazione}: Consente che un processo in modalità di sistema sia sottoposto a prelazione, rinviandone quindi
    l'esecuzione.
    \item \textbf{Kernel senza diritto di prelazione}: Non consente di applicare prelazione a processi di sistema attivi, seguitandone di conseguenza
    l'esecuzione fino a quando non cambia modalità, si blocca, oppure ceda il controllo della CPU volontariamente. Ciò rende questi kernel immuni a
    race conditions.
\end{itemize}

\noindent Si potrebbe pensare quindi che i kernel senza diritto di prelazione siano i migliori per la risoluzione del problema. Tuttavia, si tende a
preferire gli altri data la loro velocità nelle risposte e un miglior adattamento alla programmazione in tempo reale, permettendo di effettuare la
prelazione di un processo attivo nel kernel.\par
Una classica soluzione al problema della sezione critica è l'algoritmo di \textbf{Peterson}. Questo è limitato a due processi $P_i$ e $P_j$, ognuno dei
quali esegue alternativamente la propria sezione critica e non critica. È richiesto che condividano una variabile intera \textit{turn}, che indica di
chi sia il turno di accesso alla sezione, e un array di booleani \textit{flag[2]}, il quale segnala i processi pronti a entrarvi.\par
Prendendo come esempio il seguente codice del processo $P_i$, si può vedere come segnali la richiesta di entrare ponendo il proprio flag a vero e dia
il turno all'altro processo, per evitare di entrare nella sezione critica quando quest'ultimo la sta utilizzando. Dall'altra parte, al termine
dell'esecuzione in sezione critica del processo $P_j$, porrà il suo flag a falso, consentendo a $P_i$ di entrare; così via fino al termine del software.
Inoltre, in caso di richieste contemporanee verrà mantenuta solo l'ultima richiesta.
\begin{lstlisting}[language=C][h]
    // Il seguente codice vale per il processo P-i
    while (true) {
        flag[i] = true;
        turn = j;
        while (flag[j] && turn == j);
        
        // Sezione critica

        flag[i] = false;

        // Sezione non critica
    }
\end{lstlisting}

\noindent Dimostriamo adesso le tre proprietà che il protocollo deve rispettare per fornire una soluzione al problema della sezione critica: per
quanto riguarda la \textbf{mutua esclusione}, notiamo come nel codice ogni processo $P_i$ accede alla propria sezione critica se e solo se il flag di
$P_j$ è falso oppure il turn ha valore $i$; inoltre, se entrambi i processi sono eseguibili in concomitanza nelle rispettive sezioni critiche, allora
entrambi i flag saranno posti a vero. Per queste ragioni è impossibile che $P_i$ e $P_j$ eseguano le proprie istruzioni del ciclo while allo stesso
tempo. La variabile \textit{turn}, infatti, non può avere più di un valore e fino al termine della permanenza di $P_j$ nella sua sezione critica
rimangono il flag di $j$ a vero e turn con il valore di $j$, confermando la mutua esclusione.\par
Per il \textbf{progresso} e l'\textbf{attesa limitata} invece notiamo come l'ingresso di $P_i$ nella propria sezione critica è impedito solo se è
bloccato nella sua iterazione while, con variabili flag di $j$ a vero e turn con valore $j$. Non esistono altre possibilità perché: \begin{itemize}
    \item Se $P_j$ non fosse pronto ad entrare in sezione, il flag di $j$ sarebbe falso, permettendo l'accesso a $P_i$.
    \item Se il flag di $j$ è vero e il processo sta eseguendo la sua sezione critica, può avere turn con valore $i$, facendo entrare $P_i$ in critica,
    oppure turn con valore $j$, che vede $P_j$ continuare ad eseguirla.
\end{itemize}
\noindent All'uscita dalla sezione, $P_j$ imposterà il suo flag a falso per permettere a $P_i$ di entrare, mentre se $P_j$ lo reimposta a true, dovrà
anche porre turn a valore $i$. Siccome $P_i$ non modifica questi dati mentre è al di fuori della sezione critica, entrerà al suo interno dopo che $P_j$
avrà effettuato più di un ingresso, garantendo gli altri due requisiti.\par
La soluzione di Peterson posta in questi termini potrebbe tuttavia non funzionare sempre a causa di come i compilatori e i processori lavorano sulle
architetture moderne, principalmente perché attuano il riordinamento delle istruzioni di lettura e scrittura per ottimizzare le prestazioni. L'unico
modo per preservare la mutua esclusione è l'implementazione di strumenti di sincronizzazione adeguati, trattati nelle sezioni successive.

%

\section{Supporto hardware}
Con Peterson abbiamo definito una soluzione software per il problema della sezione critica, perché l'algoritmo garantisce mutua esclusione senza dover
dipendere dall'hardware. Tuttavia, questa indipendenza è ciò che di base non garantisce il funzionamento su architetture moderne, ragion per cui ci
affideremo a meccanismi di sincronizzazione.\par
Gli elaboratori utilizzano un \textbf{modello di memoria}, la modalità per determinare le garanzie sulla memoria da fornire agli applicativi. Ne
esistono di due tipi: \begin{itemize}
    \item \textbf{Fortemente ordinato}: Una modifica alla memoria su un processore è immediatamente visibile anche a tutti gli altri.
    \item \textbf{Debolmente ordinato}: Una modifica alla memoria su un processore non sempre è immediatamente visibile anche a tutti gli altri.
\end{itemize}

\noindent In quanto i modelli variano in base al processore, è necessario che forniscano istruzioni ad hoc per consentire la loro implementazione; tali
istruzioni sono dette \textbf{barriere di memoria} e alla loro esecuzione, il sistema garantisce che siano eseguite tutte le operazioni di load e
store antecedenti alle successive, risolvendo anche il problema del riordinamento delle istruzioni.\par
Le barriere di memoria fungono quindi da "blocco" per assicurarsi l'esecuzione di un insieme di operazioni prima di procedere con le successive parti
del codice; sono operazioni di basso livello spesso utilizzate nei sistemi operativi o dagli sviluppatori dei kernel.\par
Le moderne architetture forniscono anche un particolare insieme di istruzioni dedito alla manipolazione delle parole di memoria in modo
atomico: \begin{itemize}
    \item \textbf{test\_and\_set()}: Controlla e modifica il contenuto di una parola di memoria. Imposta atomicamente una variabile a true e
    restituisce il suo valore precedente, consentendo di rilevare se il lock era già acquisito.
    \item \textbf{compare\_and\_swap()}: Scambia il contenuto di due parole di memoria. Usa tre operandi: \textbf{int *value}, indicante il valore
    della parola, \textbf{int expected}, il valore che ci si aspetta abbia la parola, e \textbf{int new\_value}, il nuovo valore di quest'ultima, preso
    dall'altra parola. Se il valore della parola è uguale all'atteso, viene aggiornato con quello nuovo, altrimenti non è effettuata alcuna modifica.
    In ogni caso, ritorna il value discusso.
\end{itemize}

\noindent Queste due istruzioni sono eseguite atomicamente, quindi in presenza di due loro istanze esecutive su due unità di elaborazione, saranno
eseguite in modo sequenziale con un ordine arbitrario. Con questa dinamica e l'aiuto di una variabile booleana globale \textbf{lock}, è possibile
garantire la mutua esclusione. In particolare, si implementano il lock e un array di booleani globale inizializzati a false; quando un processo richiede
l'accesso alla propria sezione critica, setta la sua cella di array a true e dichiara una variabile key a valore 1. Fintanto che presentano questa
configurazione, ovvero fino a quando la sezione critica è utilizzata da un altro processo, il flusso rimane in un ciclo while fino a quando key non è
aggiornata con la compare and swap. Terminata l'iterazione, la cella dell'array sarà posta a false e il processo entrerà nella sua sezione critica.\par
Per quanto riguarda l'attesa limitata, invece, si itera l'array in ordine ciclico e setta il primo processo in attesa come il prossimo che può entrare
nella propria sezione critica; quindi qualunque processo voglia accedervi, potrà farlo entro $n-1$ turni.

\begin{lstlisting}[language=C]
    while (true) {
        waiting[i] = true;
        key = 1;

        while (waiting[i] && key == 1) {
            key = compare_and_swap(&lock, 0, 1);
        }
        waiting[i] = false;

        // Sezione critica

        j = (i+1) % waiting_len;
        while ((j != i) && !waiting[j]) j = (j+1) % waiting_len;

        if (j == i) lock = 0;
        else waiting[j] = false;

        // Sezione non critica
}
\end{lstlisting}
\noindent In genere, dato che sono operazioni atomiche, test\&set e compare\&swap sono utilizzate come base per costruire strumenti di sincronizzazione
più astratti, come le \textbf{variabili atomiche}, le quali forniscono operazioni su tipi di dati base, come interi e booleani. Possono essere
utilizzate per garantire mutua esclusione in alcune situazioni dove possono presentarsi race condition. Sono usate comunemente nei sistemi operativi e
applicazioni concorrenti.

%

\section{Lock Mutex, Semafori e Monitor}
Le soluzioni hardware possono essere astratte in strumenti software più accessibili come il \textbf{Lock Mutex}, usato per proteggere le sezioni
critiche e prevenire le race condition.\par
La procedura associata vede un processo acquisire il lock prima di entrare nella regione, e rilasciarlo quando ne esce. Il meccanismo è ottenuto grazie
alle funzioni \textbf{acquire()}, che acquisisce il lock quando è disponibile e pone i processi che lo richiedono in attesa attiva quando non lo è, fin
quando non viene rilasciato con \textbf{release()}. Queste devono essere eseguite atomicamente e lo stato del lucchetto è consultabile grazie a una sua
variabile booleana \textbf{available}.
\begin{lstlisting}[language=C]
    // Definizione di acquire
    acquire () {
        while (!available) { /* Attendi */}
        available = false;
    }

    // Definizione di release
    release () {
        available = true;
    }
\end{lstlisting}

\noindent Il vantaggio di questo meccanismo è di non rendere necessario alcun cambio di contesto, che tendenzialmente occupa molto tempo, tuttavia, se
molti processi richiedono di acquisire un lock indisponibile, questi continueranno a ciclare e chiamare acquire() fino al suo rilascio, sprecando cicli
della CPU. Per questa ragione sono chiamati anche "spinlock" e vengono utilizzati in sistemi multiprocessore dove possono far girare i thread su core
diversi mentre il detentore del lock esegue la sua sezione critica.\par
Uno strumento più robusto per gestire la sincronizzazione dei processi è quello dei \textbf{semafori}; sono delle variabili intere a cui si può
accedere\footnote{Escludendo l'inizializzazione.} solo tramite due operazioni atomiche chiamate \textit{wait()} e \textit{signal()}, definite come segue:
\begin{lstlisting}[language=C]
    wait (S) {
        while (S <= 0) {
            // Attesa attiva
        }
        S--;
    }

    signal (S) { S++; }
\end{lstlisting}

\noindent Tutte le modifiche al valore del semaforo $S$ devono essere effettuate all'interno di queste due funzioni e in modo indivisibile. Inoltre,
nel caso di una chiamata a \textit{wait(S)}, sarà necessario verificare anche il valore intero del semaforo e la sua eventuale modifica. Distinguiamo
due tipi di semafori: \begin{itemize}
    \item \textbf{Binari}: Il loro valore è limitato a $[0,1]$; in quanto simili ai lock mutex, vengono usati al loro posto dove i primi non possono
    essere applicati.
    \item \textbf{Contatore}: Il loro valore è un numero intero. Sono utilizzati in un contesto dove è necessario un controllo dell'accesso ad una
    risorsa presente in più esemplari.
\end{itemize}

\noindent In particolare, per questi ultimi, si inizializzano con il valore delle risorse disponibili, e i processi che ne desiderano utilizzare una
dovranno invocare \textit{wait(S)}, decrementandone quindi il valore, mentre coloro che devono restituirne, useranno \textit{signal(S)}, aumentandolo.
Nel caso in cui il semaforo valga $0$ significa che tutte le risorse sono occupate e i processi richiedenti dovranno \textbf{bloccarsi} fin quando il
semaforo non torna positivo.\par
Ciò introduce nuovamente il problema dell'attesa attiva, quindi si va a modificare le definizioni delle due operazioni elementari. Per quanto riguarda
\textit{wait()} possiamo bloccare il processo e porlo nello stato di attesa, sicché lo scheduler possa fare un cambio di contesto e dare il controllo
della CPU a un altro processo pronto. Questo processo sospeso, in attesa del semaforo, sarà riavviato in seguito ad una chiamata a \textit{signal(S)}
da parte di un altro processo e tramite un'operazione chiamata \textit{wakeup()} si modifica il suo stato da attesa a pronto, spostandolo nella relativa
coda. In un contesto tale, ridefiniamo il semaforo come segue: \begin{lstlisting}[language=C]
    typedef struct {
        int value; // Valore del semaforo
        struct process *list; // Processi in attesa
    } semaphore;
\end{lstlisting}

\noindent Mentre in pseudocodice, wait e signal diventano: \begin{lstlisting}[language=C]
    wait (semaphore *S) {
        S->value--;
        if (S->value < 0) {
            // Aggiungi il processo a S->list
            sleep();
        }
    }

    signal (semaphore *S) {
        S->value++;
        if (S->value <= 0) {
            // Rimuovi un processo da S->list
            wakeup();
        }
    }
\end{lstlisting}

\noindent La funzione \textit{sleep()} sospende il processo che la invoca, mentre \textit{wakeup()} lo pone come pronto; queste sono due operazioni
fornite dal sistema operativo tramite system call.\par
Contrariamente alla precedente definizione di semaforo, che lo mostrava come un intero, con queste modifiche è possibile scendere a numeri negativi, il
cui valore assoluto rappresenta il numero di processi in attesa che il semaforo si liberi. La lista dei processi, invece, è tipicamente implementata
con un puntatore a una coda FIFO contenente ciascun PCB.\par
Consideriamo ora le operazioni di \textit{wait()} e \textit{signal()}; in quanto è richiesto che siano atomiche, si deve garantire che nessun processo
le attui allo stesso tempo su uno stesso semaforo; ciò è riconducibile a un problema di accesso alla sezione critica. Nei sistemi a singola CPU
banalmente si inibiscono gli interrupt durante le esecuzioni delle due operazioni, ma nei sistemi multi-core disabilitare le interruzioni porterebbe a
un drastico calo delle prestazioni; per questa ragione mettono a disposizione diverse tecniche di realizzazione dei lock, come spinlock o
\textit{compare\_and\_swap()}.\par
Notiamo infine che la nuova definizione di \textit{wait()} e \textit{signal()} non elimina completamente il problema dell'attesa attiva, bensì lo
rimuove dalle sezioni di ingresso degli applicativi; poiché infatti questa si limita alle sezioni critiche delle due operazioni, queste sono in genere
convenientemente codificate, sicché siano particolarmente brevi. Ciò fa presentare raramente l'attesa attiva e quando accade, lo fa per breve
tempo.\newline 

\noindent Nonostante i semafori siano un costrutto pratico ed efficace, un loro uso scorretto può causare vari problemi di difficile individuazione,
in quanto si possono manifestare solo in presenza di particolari sequenze di esecuzione. Non escludono nemmeno il verificarsi di problemi di
sincronizzazione; supponiamo infatti il problema della sezione critica dove tutti i processi condividono una variabile semaforo mutex inizializzata a
$1$. Ogni processo chiama \textit{wait(mutex)} all'entrata nella propria sezione critica e invoca \textit{signal(mutex)} quando ne esce. Se la sequenza
non è rispettata possono verificarsi i seguenti problemi: \begin{itemize}
    \item \textbf{Ordine della operazioni inverso}: Più processi possono eseguire le proprie sezioni critiche, violando la mutua esclusione.
    \item \textbf{Sostituzione di signal(mutex) con wait(mutex)}: Il processo si bloccherà permanentemente alla seconda chiamata \textit{wait()},
    perché il semaforo non sarà più disponibile.
    \item \textbf{Omissione di una o entrambe le operazioni}: Si viola la mutua esclusione oppure il processo si blocca indefinitamente.
\end{itemize}

\noindent Una strategia per la gestione di questi errori consiste nell'incorporare ulteriori strumenti di sincronizzazione astratti, come il tipo
\textbf{monitor}.\par
Un monitor è un \textbf{tipo di dato astratto} (ADT) che incapsula i dati mettendo a disposizione un insieme di funzioni, contraddistinte dalla mutua
esclusione, per operare su di essi. Contiene inoltre la dichiarazione delle variabili i cui valori definiscono lo stato dell'istanza del monitor.
Generalmente, la sintassi di un monitor in pseudocodice è: \begin{lstlisting}[language=C]
    monitor monitorName {
        // Dichiarazione di variabili condivise

        function P1 (...) { ... }
        function P2 (...) { ... }

        // Altre funzioni

        function Pn (...) { ... }
        initCode (...) { ... }
    }
\end{lstlisting}

\noindent Le funzioni dentro al monitor hanno accesso solo alle variabili locali o ai parametri formali. Il costrutto assicura inoltre che al suo
interno possa essere attivo un solo processo alla volta grazie ad eventuali variabili \textbf{condition} $x$, costrutti sui quali sono applicabili
esclusivamente le operazioni di \textit{wait()} e \textit{signal()}.\par
Entrambe mantengono la stessa funzione discussa con i semafori, ma la seconda non ha effetto in assenza di processi sospesi. Supponiamo che un processo
$P$ chiami \textit{x.signal()} e che ne esista un altro $Q$ associato alla condition $x$. Se si permette a $Q$ di riprendere la propria esecuzione,
l'altro processo $P$ dovrà necessariamente attendere. Tuttavia, concettualmente, potrebbero entrambi continuare l'esecuzione, creando due
possibilità: \begin{itemize}
    \item \textbf{Segnalare e attendere}: $P$ attende che $Q$ lasci il monitor o attenda su un'altra condition.
    \item \textbf{Segnalare e continuare}: $Q$ attende che $P$ lasci il monitor o attenda su un'altra condition.
\end{itemize}

\noindent Non esiste un metodo migliore rispetto all'altro, ma per le implementazioni useremo la prima istanza. Presentiamo una possibile realizzazione
di un monitor tramite semafori. A ogni monitor sarà associato un semaforo mutex inizializzato a $1$ per garantire la mutua esclusione. Un processo
eseguirà \textit{wait(mutex)} prima di entrare nel meccanismo e chiamerà \textit{signal(mutex)} quando uscirà dal monitor. In quest'ultimo scenario,
dato che il processo chiamante dovrà attendere finché quello risvegliato si metta in attesa o lasci il monitor, si introduce un nuovo semaforo
\textit{next} inizializzato a $0$, su cui i processi segnalanti possono autosospendersi. Il relativo contatore si implementa con una variabile intera
\textit{next\_count}. Si richiede quindi che ogni procedura esterna $F$ sia compresa nel seguente codice: \begin{lstlisting}[language=C]
    wait(mutex);

    // Corpo di F

    if (next_count > 0) signal(next);
    else signal(mutex);
\end{lstlisting}

\noindent Ciò assicura la mutua esclusione dentro al monitor. La realizzazione delle condition vede poi l'introduzione di un semaforo \textit{x\_sem} e
un intero \textit{x\_count}, ambo inizializzati a $0$. Quindi, le operazioni di \textit{wait()} e \textit{signal()} associate alla condition $x$ si
ridefiniscono come segue: \begin{lstlisting}[language=C]
    wait () {
        x_count++;
        if (next_count > 0) signal(next);
        else signal(mutex);
        wait(x_sem); // Chiamata a wait di semaforo, non condition
        x_count--;
    }

    signal () {
        if (x_count > 0) {
            next count++;
            signal(x_sem); // Chiamata a signal di semaforo, non condition
            wait(next); // Idem qua
            next_count--;
        }
    }
\end{lstlisting}

\noindent È necessario ora decidere l'ordine di ripresa dei processi sospesi alla condition $x$. Generalmente si crea un costrutto di
\textbf{attesa condizionale} con forma \textit{x.wait(c)}, dove $c$ è un valore di \textbf{priorità} determinato da un'espressione valutata al momento
dell'esecuzione di \textit{wait()}. Tale numero è poi associato a un processo sospeso, e all'esecuzione di \textit{x.signal()} si riprenderà dal
processo con la priorità più alta in base alla convenzione scelta.\par
Di base, tuttavia, il monitor non garantisce il rispetto di questo ordine, poiché potrebbero accadere i seguenti scenari: \begin{itemize}
    \item Un processo accede alla risorsa senza prima ottenere il permesso.
    \item Ottenuto l'accesso alla risorsa, un processo potrebbe non rilasciarla più.
    \item Un processo tenta di rilasciare una risorsa che non ha mai richiesto.
    \item Un processo può richiedere più volte la stessa risorsa senza averla rilasciata prima delle successive richieste.
\end{itemize}

\noindent È evidente che bisogna concentrarsi su eventuali soluzioni in merito alla gestione delle operazioni ad alto livello, controllando i programmi
che utilizzano il monitor in tal modo che si verifichino due condizioni: in primo luogo, i \textbf{processi utenti} devono sempre impiegare il costrutto
seguendo la \textbf{sequenza corretta}; in secundis, bisognerà assicurare che eventuali \textbf{processi non cooperativi} non \textbf{aggirino} la
\textbf{mutua esclusione} offerta.\par
Rispettando le due condizioni si ha la certezza di non avere problemi di sincronizzazione, ma è applicabile agevolmente solo per sistemi di piccole
dimensioni.

%

\section{Liveness}
L'impiego di strumenti di sincronizzazione può avere conseguenze relative all'attesa indefinita a determinate risorse. Con \textbf{liveness} intendiamo
un insieme di proprietà che un sistema deve soddisfare per evitare questo problema e garantire che ogni processo progredisca la sua esecuzione durante
il proprio ciclo vitale. Strumenti come i lock mutex e semafori consentono di indurre i problemi di \textbf{mancanza di liveness} alla programmazione
concorrente. Analizziamo ora alcune istanze di questo tipo di problemi con relative soluzioni.\par
Il primo caso di studio è rappresentato dallo \textbf{stallo} (deadlock); una situazione nella quale uno o più processi rimangono in eterna
attesa del verificarsi di un evento. In particolare, nelle implementazioni dei semafori con coda di attesa, i processi aspetteranno una chiamata a
\textit{signal()} impossibile da attuare.\par
Supponiamo due processi $P_0, P_1$, ognuno dei quali con un accesso a due semafori $S, Q$, ambo impostati al valore $1$. Se i due processi chiamano
\textit{wait(S)} e \textit{wait(Q)} rispettivamente allo stesso tempo, allora entrambi si bloccheranno fino alla ricezione di \textit{signal()}, ma
essendo in attesa, non possono effettuare chiamate a funzioni, risultando in stallo. La soluzione è fornita con un opportuno utilizzo della
sincronizzazione.\par
Un'altra istanza di liveness failure si può ottenere con lo scheduling dei processi. Se è implementato un meccanismo basato sulla priorità, è probabile
che un processo a priorità alta richieda dei dati utilizzati da un altro con priorità minore, ed essendo i dati a livello kernel protetti da un lock,
il primo processo sarà costretto ad attendere fino al rilascio delle suddette risorse.\par
La soluzione a questo problema, detto \textbf{inversione della priorità}, si ha implementando un protocollo per l'\textbf{ereditazione delle priorità}.
Fondamentalmente, quando un processo a bassa priorità utilizza dei dati che possono essere richiesti da altri processi più importanti, eredita il
livello di priorità di questi ultimi per tutto il tempo in cui usufruiscono di queste risorse, impedendo che processi a priorità intermedia sottraggano
la CPU al processo a bassa priorità che detiene il lock, riducendo così il tempo di attesa del processo ad alta priorità.